
This work is positioned at the intersection of reproducibility assessment, benchmark methodology, and dataset documentation standards.

\subsubsection{Reproducibility Crisis and Assessment}

The machine learning reproducibility crisis has received systematic attention. Kapoor and Narayanan (2022) developed a taxonomy of eight data leakage types, documenting their prevalence across 329 papers in 17 scientific fields. Hullman et al. (2022) drew parallels between psychology's replication crisis and machine learning, highlighting shared patterns of selective reporting and specification gaming.

Assessment tools emerged in response. rliable \citep{agarwal2021deep} provides statistical methods for reliable evaluation on reinforcement learning benchmarks. Reproscreener \citep{bhaskar2024reproscreener} leverages language models to automatically assess computational reproducibility of ML pipelines, producing a ReproScore metric. paper-replay enables end-to-end verification through an init-setup-verify-attest workflow with GPG signing.

These tools operate post-hoc. The present work differs by examining prediction of variance from dataset metadata before experimental execution.

\subsubsection{Benchmark Methodology and Standards}

The ML benchmarking community has developed infrastructure for standardized evaluation. OpenML \citep{vanschoren2014openml} provides a platform with extensive metadata including 38 auto-computed meta-features per dataset. MLCommons established industry benchmarks with specified configurations. HPOBench \citep{eggensperger2021hpobench} uses containerization to ensure environmental reproducibility.

OpenML's benchmark suites offer curated collections with machine-readable metadata. The AutoML Benchmark \citep{gijsbers2019automl} enables reproducible AutoML system evaluation across standardized datasets. These efforts focus on infrastructure reproducibility---ensuring the same code produces the same results---rather than predicting which datasets will exhibit low variance across independent implementations.

\subsubsection{Documentation and Metadata Quality}

Dataset documentation has received increasing attention. Datasheets for Datasets \citep{gebru2021datasheets} proposed structured documentation covering motivation, composition, collection process, and intended uses. Model cards \citep{mitchell2019model} provide analogous documentation for ML models. HuggingFace dataset cards operationalize these principles.

Prior work treats documentation as a qualitative good practice rather than quantifying its relationship to downstream reproducibility. This work differs by examining whether metadata completeness---operationalized as a 5-field checklist---is associated with measurable reductions in reproducibility variance.

\subsubsection{Variance Decomposition in ML}

Bouthillier et al. (2021) decomposed variance in deep learning experiments, identifying sources including random initialization, data shuffling, and hardware non-determinism. Picard (2021) quantified variation due to random seed selection. These studies characterize \textit{sources} of variance but do not predict which experimental configurations will exhibit high variance.

This work complements variance decomposition by identifying an \textit{upstream predictor}---metadata completeness---that may forecast variance magnitude before experiments run.

